%-------------------------------%
%  Author: Alessandro Sciarra   %
%    Date: 12 Jun 2025          %
%-------------------------------%

\AddLinkToTOCfalse
%~~~~~~~~~~~~~~~~~~~~~~~~~~~~~~~~~~~~~~~~~~~~%
\begin{frame}{Training schedule}
    \begin{center}
        \includegraphics[width=0.8\textwidth]{Timetable}
    \end{center}
\end{frame}
%~~~~~~~~~~~~~~~~~~~~~~~~~~~~~~~~~~~~~~~~~~~~%
\AddLinkToTOCtrue
\begin{frame}{Our mental setup}
    \begin{itemize}
        \item \PS{Relaxed and respectful atmosphere}\\ \then We are all here to learn!
        \item \textbf{It will be intense!}
        \item More practice than theory $\to$ Give your best!
        \item It's going to be interactive $\to$ ask, ask, ask!
        \item \PP{This is a changing point, but yet only a starting one!}
    \end{itemize}
    \vspace{-1mm}
    \begin{varblock*}{alert}[0.55\textwidth]{Disclaimer}
        Slides are often full of text\\
        \then They are meant for offline reading!
    \end{varblock*}
    \begin{tikzpicture}[remember picture, overlay]
        \node[anchor=north east, inner sep=5mm] at (current page.north east)
            {\includegraphics[width=0.35\textwidth]{Relax}};
        \node[anchor=south, rounded corners=3pt, draw, PS, font=\footnotesize] (L) at ($(current page.south)+(0,0.6)$)
            {Clickable link to clickable ToC};
        \draw[to] (L.south) |- ($(current page.south east)+(-0.5,0.25)$);
    \end{tikzpicture}
\end{frame}
\AddLinkToTOCfalse
%~~~~~~~~~~~~~~~~~~~~~~~~~~~~~~~~~~~~~~~~~~~~%
\begin{frame}{Alessandro -- Who am I?}
    \setbeamertemplate{itemize/enumerate body begin}{\small}
    \setbeamertemplate{itemize/enumerate subbody begin}{\footnotesize}
    \begin{varblock}{quote}[0.75\textwidth]{Ancient popular wisdom}
        \small\guillemotleft{}I am not saying to be Superman, I am just saying nobody has ever seen me and Superman together in the same room.\guillemotright
    \end{varblock}
    \vspace{5mm}
    \begin{itemize}
        \item Member of the Software Development Center \\ $\to\;$ Z02 project of the CRC-TR\,211 collaboration
        \item A used-to-be ``lattice practitioner''
        \item A dreamer who flatters himself that
              \begin{itemize}
                  \footnotesize
                  \item physicists will one day \PB{properly} use the IT tools they use
                  \item IT people will teach IT to physicists and work next to them
                  \item physicists will be able to focus on physics
              \end{itemize}
        \item \PP{Not really an expert, but rather a self-thought software enthusiast}
    \end{itemize}
    \begin{tikzpicture}[remember picture, overlay]
        \node[anchor=south east, inner sep=0] at (current page.south east) {
\includegraphics[height=4cm]{Superman}};
    \end{tikzpicture}
\end{frame}
%~~~~~~~~~~~~~~~~~~~~~~~~~~~~~~~~~~~~~~~~~~~~%
\begin{frame}{David -- Who am I?}
\begin{itemize}
    \item Postdoc at Bielefeld
    \item Lattice practitioner studying QCD phase diagram,
          muon $g-2$
    \item Lots of experience developing high-performance
          lattice QCD code
          \begin{itemize}
              \item \URL*{https://github.com/LatticeQCD/SIMULATeQCD}{SIMULATeQCD}
              \item \URL*{https://github.com/paboyle/Grid}{Grid}
          \end{itemize}
    \item Data analysis: Collect generic tools
          in \URL*{https://github.com/LatticeQCD/AnalysisToolbox}{AnalysisToolbox}
    \item Not really expert either; rather physicist
          who is paranoid about being wrong and perhaps guilty
          of improperly using IT tools sometimes :)
\end{itemize}
\end{frame}
%~~~~~~~~~~~~~~~~~~~~~~~~~~~~~~~~~~~~~~~~~~~~%
\begin{frame}{What to expect from this training}
    \vspace{0mm}
    \begin{columns}
        \hfill
        \begin{column}{0.38\textwidth}
            \begin{center}
                \includegraphics[width=0.9\linewidth]{Old-new-way}
            \end{center}
        \end{column}
        \begin{column}{0.6\textwidth}
            \begin{varblock}{quote}[\linewidth]{\small A scientist's point of view}
                \footnotesize
                The way you program today is better than it was yesterday and probably worse than it will be tomorrow!
                Have a good reason to do things the way you do them.
                It will happen that you will change mind and consequently your approach to the same problem.
                Again and over again, it is normal.
                Do not hesitate in doing so.
                To change mind and approach means to evolve and to become a better programmer.
            \end{varblock}
        \end{column}
        \hfill
    \end{columns}
    \vspace{3mm}
    \begin{center}
        {\Large\PS{Embrace a new working philosophy!}}\\[1mm]
        $\downarrow$\\[1mm]
        \PB{High-value material if you ever quit academia}
    \end{center}
\end{frame}
%~~~~~~~~~~~~~~~~~~~~~~~~~~~~~~~~~~~~~~~~~~~~%
\begin{frame}{A word about programming languages}
    \begin{varblock}{alert}[0.95\textwidth]{How can we work and learn together if we use different languages?}
        \begin{itemize}
            \item Most principles are language independent
            \item All used language in science are somewhat similar
            \item Trust your colleagues and focus on ideas, not technicalities
        \end{itemize}
    \end{varblock}
    \begin{varblock}{quote}[\textwidth]{The Pragmatic Programmer}[A.~Hunt \& D.~Thomas]
        «Learn at least one new language every year.
        Different languages solve the same problems in different ways.
        By learning several different approaches, you can help broaden your thinking and avoid getting struck in a rut.»
    \end{varblock}
    \begin{center}
        \Large
        \PS{Be excited, never frustrated!}
    \end{center}
\end{frame}
%~~~~~~~~~~~~~~~~~~~~~~~~~~~~~~~~~~~~~~~~~~~~%
\begin{frame}{What do we expect from this training?}
    \begin{itemize}
        \item We are aware that a lot will might be new
        \item \PS{Do not hesitate to ask!}
        \item Speak up in case of need $\;\to\;$ \alert{It should never get frustrating!}
        \item \PP{Put effort into the exercises}
        \item The more active you are the more information will stay in your memory
    \end{itemize}
    \vspace{3mm}
    \begin{center}
        \includegraphics[height=0.36\textheight]{Top}
        
\includegraphics[height=0.36\textheight]{Tired-but-happy}
    \end{center}
    \begin{tikzpicture}[remember picture, overlay, scope on=<2>]
        \node (H) at ([yshift=5mm]current page.center) {
\includegraphics[height=0.6\textheight]{Excited}};
        \node[below = 1mm of H, font=\Huge\bfseries, text=PT, fill=BGLIGHT, inner sep=1mm] {Are you ready?};
    \end{tikzpicture}
\end{frame}
%~~~~~~~~~~~~~~~~~~~~~~~~~~~~~~~~~~~~~~~~~~~~%
\AddLinkToTOCtrue
